% Source context: GLM23, arXiv:2203.12563v3, REsubmission.tex,
% Section 5, lines 1695--1777. The coefficient-one open-chain penalty
% below is an auxiliary consequence of the commuting constraints.
% Native declaration and diagram validation is pending for this source leaf.

\subsection{The active range of the actual joint open chain}
\label{sec:mpo_joint_open_sector_penalty}

Write the original physical space as
$\mathcal H=\mathcal H_{00}\oplus\mathcal H_{01}
    \oplus\mathcal H_{10}\oplus\mathcal H_{11}$, where
$\mathcal H_{00}=\C^{d_0}$ is shared by all block labels.
Let $R_0$ project onto $\mathcal H_{00}\oplus\mathcal H_{01}$,
let $C_0$ project onto $\mathcal H_{00}\oplus\mathcal H_{10}$,
and set $P_{00}=R_0C_0=C_0R_0$. Let $S$ be the actual zero-parameter
two-site support and $h=I-P_S$ its complementary projection.
For the full joint polar boundary frames define
$E_L=P_{\operatorname{ran}U_L}$ and
$E_R=P_{\operatorname{ran}U_R}$.

\begin{definition}[The inner phase constraint]
    \label{def:mpo_joint_inner_projection}
    \lean{MPSTensor.MPOSymmetry.jointMixedTwoSiteInnerProjection}
    The inner phase constraint on two sites is $F_{\mathrm{in}}=C_0\otimes R_0$.
    It selects column phase zero at the first site and row phase zero
    at the second site.
\end{definition}

\begin{theorem}[The actual support satisfies the inner constraint]
    \label{thm:mpo_joint_inner_projection_support}
    \lean{MPSTensor.MPOSymmetry.jointMixedTwoSiteInnerProjection_isSymmetricProjection,
        MPSTensor.MPOSymmetry.range_blockInsertedBoundaryMap_jointMixed_le_innerProjection,
        MPSTensor.MPOSymmetry.one_sub_jointMixedTwoSiteInnerProjection_le_parentInteraction}
    \uses{def:mpo_joint_inner_projection}
    The operator $F_{\mathrm{in}}$ is an orthogonal projection, and
    $S\subseteq\operatorname{ran}F_{\mathrm{in}}$. Consequently
    $I-F_{\mathrm{in}}\le h$.
\end{theorem}
\begin{proof}
    The two commuting zero-one diagonal selectors are orthogonal.
    For every actual support vector $v$, the inserted trace contraction
    gives $(C_0\otimes I)v=v=(I\otimes R_0)v$. Their product therefore
    fixes $v$. The inclusion implies $P_S\le F_{\mathrm{in}}$ and
    hence $I-F_{\mathrm{in}}\le I-P_S=h$.
\end{proof}

\begin{theorem}[Phase selectors reduce all translated interactions]
    \label{thm:mpo_joint_open_phase_commutators}
    \lean{MPSTensor.MPOSymmetry.jointMixedRowSector_commute_periodicLocalInteraction_zero,
        MPSTensor.MPOSymmetry.jointMixedColumnSector_commute_periodicLocalInteraction_zero,
        MPSTensor.MPOSymmetry.jointMixedRowSector_commute_openInteractionHamiltonian_zero,
        MPSTensor.MPOSymmetry.jointMixedColumnSector_commute_openInteractionHamiltonian_zero}
    Let $N\ge2$. Write $h_i$ for the cyclic translate of $h$ to sites
    $i,i+1$ and $H_N=\sum_{i=0}^{N-2}h_i$ for the nonwrapping open sum.
    For all chain sites $k$ and cyclic starts $i$,
    \begin{align}
        [(R_0)_k,h_i] & =0, & [(C_0)_k,h_i] & =0, &
        [(R_0)_k,H_N] & =0, & [(C_0)_k,H_N] & =0.
    \end{align}
\end{theorem}
\begin{proof}
    Both local row selectors and both local column selectors commute
    with $h$: the inner selectors fix $S$, while the outer selectors
    preserve $S$ by acting on its virtual boundary. A selector inside
    a translated window therefore commutes by its local identity.
    A selector outside the window is constant on each active-coordinate
    fiber and also commutes. Summing the nonwrapping translates gives
    the two commutators with $H_N$.
\end{proof}

\begin{definition}[The open active projection]
    \label{def:mpo_joint_open_active_projection}
    \lean{MPSTensor.MPOSymmetry.jointMixedOpenLeftMatrix,
        MPSTensor.MPOSymmetry.jointMixedOpenRightMatrix,
        MPSTensor.MPOSymmetry.jointMixedOpenBondProjection,
        MPSTensor.MPOSymmetry.jointMixedOpenActiveProjection}
    For $N\ge3$, write $h_i$ for $h$ on sites $i,i+1$ and the
    identity elsewhere, and set $H_N=\sum_{i=0}^{N-2}h_i$.
    Define projections on the same bonds by
    \begin{align}
        F_0     & =(E_L)_0(R_0)_1,                         &
        F_i     & =(C_0)_i(R_0)_{i+1}\quad(1\le i\le N-3), &
        F_{N-2} & =(C_0)_{N-2}(E_R)_{N-1}.
        \label{eq:mpo_joint_open_bond_constraints}
    \end{align}
    Their product is
    \begin{align}
        Q_N & =F_0F_1\cdots F_{N-2}.
        \label{eq:mpo_joint_open_active_product}
    \end{align}
    The middle family is empty at $N=3$.
\end{definition}

\begin{theorem}[The exact open active range]
    \label{thm:mpo_joint_open_active_range}
    \lean{MPSTensor.MPOSymmetry.jointMixedOpenBondProjection_commute,
        MPSTensor.MPOSymmetry.jointMixedOpenBondProjection_isSymmetricProjection,
        MPSTensor.MPOSymmetry.mem_range_jointMixedOpenActiveProjection_iff_bonds,
        MPSTensor.MPOSymmetry.jointMixedOpenActiveProjection_isSymmetricProjection,
        MPSTensor.MPOSymmetry.mem_range_jointMixedOpenActiveProjection_iff}
    \uses{def:mpo_joint_open_active_projection}
    The $F_i$ commute, and each $F_i$ and $Q_N$ is an orthogonal projection.
    The range of $Q_N$ is the common fixed space of both one-site factors
    in every $F_i$. Equivalently, a vector $v$ belongs to its range precisely when
    \begin{align}
        (E_L)_0v & =v,                              & (E_R)_{N-1}v & =v, &
        (R_0)_kv & =v=(C_0)_kv\quad(1\le k\le N-2).
        \label{eq:mpo_joint_open_active_conditions}
    \end{align}
    Thus the active subspace is
    $\operatorname{ran}U_L\otimes\mathcal H_{00}^{\otimes(N-2)}
        \otimes\operatorname{ran}U_R$.
    Every original direction of $\mathcal H_{00}$ is retained,
    including directions unused by the endpoint tensor.
\end{theorem}
\begin{proof}
    A boundary frame occurs only at its boundary site. Distinct bonds
    therefore overlap only in $R_0$ and $C_0$, which commute.
    The product of commuting orthogonal projections is orthogonal and
    fixes exactly the vectors fixed by each factor.
    Applied to (\ref{eq:mpo_joint_open_bond_constraints}), this is
    (\ref{eq:mpo_joint_open_active_conditions}). At an interior site the
    two phase-zero conditions select exactly the shared $00$ summand.
\end{proof}

For four sites, the product in
(\ref{eq:mpo_joint_open_active_product}) has the following contraction.
Each vertical line is the full original physical space $\mathcal H$.
\begin{tenkzequation}
    % Formula: (E_L tensor R_0 tensor I tensor I)
    % (I tensor C_0 tensor R_0 tensor I)
    % (I tensor I tensor C_0 tensor E_R)
    % = E_L tensor P_00 tensor P_00 tensor E_R.
    % Boundary signature on both sides: up=4, down=4, west=0, east=0.
    % The rows act from bottom input to top output. All wires have dimension dim H.
    \tenkzkernel
    \begin{tenkz}[rows={wire,wire,wire}, cols=4, bonds=none]
        \tn[at=(1,1), name=A, skin=box, size=m,
            ports={90:physical:$i_0$,270:physical}]{E_L}
        \tn[at=(1,2), name=B, skin=box, size=m,
            ports={90:physical:$i_1$,270:physical}]{R_0}
        \tn[at=(1,3), name=C, skin=box, size=m,
            ports={90:physical:$i_2$,270:physical}]{I}
        \tn[at=(1,4), name=D, skin=box, size=m,
            ports={90:physical:$i_3$,270:physical}]{I}
        \tn[at=(2,1), name=E, skin=box, size=m,
            ports={90:physical,270:physical}]{I}
        \tn[at=(2,2), name=F, skin=box, size=m,
            ports={90:physical,270:physical}]{C_0}
        \tn[at=(2,3), name=G, skin=box, size=m,
            ports={90:physical,270:physical}]{R_0}
        \tn[at=(2,4), name=H, skin=box, size=m,
            ports={90:physical,270:physical}]{I}
        \tn[at=(3,1), name=J, skin=box, size=m,
            ports={90:physical,270:physical:$j_0$}]{I}
        \tn[at=(3,2), name=K, skin=box, size=m,
            ports={90:physical,270:physical:$j_1$}]{I}
        \tn[at=(3,3), name=L, skin=box, size=m,
            ports={90:physical,270:physical:$j_2$}]{C_0}
        \tn[at=(3,4), name=M, skin=box, size=m,
            ports={90:physical,270:physical:$j_3$}]{E_R}
        \tnwire{A.270}{E.90}
        \tnwire{B.270}{F.90}
        \tnwire{C.270}{G.90}
        \tnwire{D.270}{H.90}
        \tnwire{E.270}{J.90}
        \tnwire{F.270}{K.90}
        \tnwire{G.270}{L.90}
        \tnwire{H.270}{M.90}
    \end{tenkz}
    \quad $=$ \quad
    \begin{tenkz}[rows={wire}, cols=4, bonds=none]
        \tn[at=(1,1), name=A, skin=box, size=m,
            ports={90:physical:$i_0$,270:physical:$j_0$}]{E_L}
        \tn[at=(1,2), name=B, skin=box, size=m,
            ports={90:physical:$i_1$,270:physical:$j_1$}]{P_{00}}
        \tn[at=(1,3), name=C, skin=box, size=m,
            ports={90:physical:$i_2$,270:physical:$j_2$}]{P_{00}}
        \tn[at=(1,4), name=D, skin=box, size=m,
            ports={90:physical:$i_3$,270:physical:$j_3$}]{E_R}
    \end{tenkz}
\end{tenkzequation}

\begin{theorem}[Each active constraint is controlled by its own bond]
    \label{thm:mpo_joint_open_bond_control}
    \lean{MPSTensor.MPOSymmetry.one_sub_jointMixedOpenBondProjection_le_localInteraction,
        MPSTensor.MPOSymmetry.jointMixedOpenBondProjection_commute_localInteraction}
    \uses{def:mpo_joint_open_active_projection}
    For every pair of nonwrapping bonds $i,j$,
    $I-F_i\le h_i$ and $[F_i,h_j]=0$.
\end{theorem}
\begin{proof}
    \uses{thm:mpo_joint_one_sided_support,thm:mpo_joint_inner_projection_support,
        thm:mpo_joint_open_phase_commutators}
    The first and last bond inequalities are the two one-sided support
    inequalities; the remaining bonds use $I-F_{\mathrm{in}}\le h$.
    Extension by the spectator identity preserves order.
    Every phase factor of $F_i$ commutes with $h_j$ by the phase
    commutators. Its possible first or last frame factor commutes with
    its incident interaction by support inclusion and with every other
    interaction by disjoint support. Thus both factors, and their
    product $F_i$, commute with $h_j$.
\end{proof}

\begin{theorem}[Reduction and the inactive-sector penalty]
    \label{thm:mpo_joint_open_inactive_penalty}
    \lean{MPSTensor.MPOSymmetry.jointMixedOpenActiveProjection_commute_openInteraction,
        MPSTensor.MPOSymmetry.one_sub_jointMixedOpenActiveProjection_le_openInteraction,
        MPSTensor.MPOSymmetry.one_sub_jointMixedOpenActiveProjection_le_twice_openInteraction}
    \uses{def:mpo_joint_open_active_projection}
    The actual open Hamiltonian satisfies
    \begin{align}
        [Q_N,H_N] & =0, & I-Q_N & \le H_N\le2H_N.
        \label{eq:mpo_joint_open_inactive_penalty}
    \end{align}
    The coefficient-one estimate is an auxiliary consequence of the
    chosen commuting constraints. It makes no assertion about the gap
    of $H_N$ within $\operatorname{ran}Q_N$.
\end{theorem}
\begin{proof}
    \uses{thm:mpo_joint_open_bond_control,thm:mpo_joint_open_active_range}
    The actual support is fixed by each local constraint, so
    $I-F_i\le h_i$. The row and column phase selectors commute with
    every translated interaction. The first frame commutes with $h_0$
    by its support inclusion and with the remaining terms by disjoint
    support; the last frame has the reflected property.
    Hence $[F_i,h_j]=0$ for every pair of nonwrapping bonds, proving
    $[Q_N,H_N]=0$. The orthogonal projection inequality
    $I-PQ\le(I-P)+(I-Q)$ for commuting $P,Q$ gives by induction
    \begin{align}
        I-Q_N & \le\sum_{i=0}^{N-2}(I-F_i)
        \le\sum_{i=0}^{N-2}h_i=H_N.
    \end{align}
    Finally $H_N\ge0$, which yields the coefficient-two bound.
\end{proof}

\begin{theorem}[The single two-site boundary term]
    \label{thm:mpo_joint_open_two_site_reduction}
    \lean{MPSTensor.MPOSymmetry.jointMixedTwoSiteFrameProjection_openInteraction_reduction}
    At $N=2$, set $Q_2=P_{\operatorname{ran}(U_L\otimes U_R)}$.
    The open sum contains one term, $H_2=h$, and satisfies
    $[Q_2,H_2]=0$ and $I-Q_2\le H_2$.
\end{theorem}
\begin{proof}
    \uses{thm:mpo_joint_full_frame_reduction}
    There is exactly one nonwrapping two-site window. Its full-frame
    support inclusion supplies both conclusions, with both boundary
    frames acting on this same term.
\end{proof}
